\section{Technology Stack and Software Artifacts Inventory}
\label{sec:impl_stack}
The EffectGuard framework and AgentSentinel evaluation testbed are implemented in modern Python 3.12 using standard, robust software engineering practices. The system is packaged using the \texttt{uv} packaging engine and configured via \texttt{pyproject.toml}. Table~\ref{tab:impl_artifacts} provides an exhaustive inventory of the primary software modules and their operational roles across the codebase.

\begin{table}[H]
\centering
\small
\renewcommand{\arraystretch}{1.2}
\begin{tabularx}{\textwidth}{p{3.8cm}p{4.2cm}X}
\toprule
\textbf{Module Path} & \textbf{Component Name} & \textbf{Core Implementation Responsibility} \\
\midrule
\texttt{agent/agent.py} & Agent Orchestrator & LangGraph state graph, Groq MoE model binding, structured output loop. \\
\texttt{agent/tools.py} & Tool Definitions & Tool schemas, date calculations, item staging, and EffectGuard gateway wrappers. \\
\texttt{agent/db.py} & Relational State Layer & Persistent SQLite connection pooling, WAL mode, item status updates, username normalization. \\
\texttt{agent/fault\_injector.py} & Chaos Fault Engine & Simulates distributed timeouts (504) pre- and post-commit, server errors (500). \\
\texttt{effectguard/gateway.py} & Execution Gateway & SHA-256 idempotency keying, WAL journaling, Verify-Before-Retry timeout recovery. \\
\texttt{effectguard/policy\_engine.py}& Policy-as-Code Engine & Invariant validation (user existence, positive amounts, safety ceilings, balance caps). \\
\texttt{effectguard/claim\_auditor.py}& Claim-vs-Ledger Auditor & Mathematical diffing of agent assertions against SQLite tickets (False Claims, Ghost Mutations). \\
\texttt{effectguard/condition\_runner.py}& Matrix Ablation Suite & Head-to-head empirical testing across conditions C0 through C4 under chaos. \\
\texttt{evaluations/dataset.py} & Golden Test Dataset & 25 rigorous test scenarios categorized across 5 enterprise operational domains. \\
\texttt{backend/eval\_service.py} & Telemetry Streaming Engine & Asynchronous SSE streaming runner computing LCS sequence metrics and financial safety. \\
\bottomrule
\end{tabularx}
\caption{Software artifacts inventory and architectural responsibilities.}
\label{tab:impl_artifacts}
\end{table}

\section{Agent Runtime and Tool Execution Pipeline}
\label{sec:impl_agent}
The agent runtime is configured in \texttt{agent/agent.py} utilizing the LangChain ReAct agent builder connected to the Groq API:
\begin{lstlisting}[language=Python]
from langchain.agents import create_agent
from langchain.agents.structured_output import ToolStrategy
from langgraph.checkpoint.memory import InMemorySaver
from models import model
from prompts import agent_system_prompt
from objects import Answer
from tools import calculate_right_refund, get_user_purchase_history, get_today_date, issue_refund_ticket

agent = create_agent(
    model=model,
    system_prompt=agent_system_prompt,
    tools=[calculate_right_refund, get_user_purchase_history, get_today_date, issue_refund_ticket],
    response_format=ToolStrategy(Answer),
    checkpointer=InMemorySaver(),
)
\end{lstlisting}

\subsection{Tool Wrapping and Staged Parameter Staging}
In \texttt{agent/tools.py}, tool implementations enforce business validation while routing state-mutating actions through EffectGuard. A critical engineering challenge arises in multi-item refund calculations: when a customer requests returns for multiple items, \texttt{calculate\_right_refund} inspects individual purchase dates against the 30-day cutoff, identifies eligible items, and computes the total amount.

To ensure that \texttt{issue\_refund\_ticket} permanently updates \textit{only} the verified eligible items in SQLite, we implemented an in-memory staging mechanism:
\begin{lstlisting}[language=Python]
# In agent/tools.py: stage verified eligible items
if eligible_items:
    stage_pending_refund_items(norm_user, eligible_items, total_refund)
\end{lstlisting}
When \texttt{issue\_refund\_ticket} subsequently commits the transaction, the staged item names are consumed atomically to flag their database status as \texttt{'REFUNDED'}.

\section{EffectGuard Gateway Core Implementation}
\label{sec:impl_gateway}
The core gateway is implemented in \texttt{effectguard/gateway.py}. The execution pipeline intercepts mutating tools, executes deterministic gating, records write-ahead logs, and manages distributed recovery:

\begin{lstlisting}[language=Python]
def execute(self, tool_name: str, args: Dict[str, Any], raw_executor: Callable, thread_id: Optional[str] = None) -> str:
    if not self.enabled:
        return raw_executor(args, "") # Naive passthrough mode (Condition C0)

    effect_class = policy_engine.get_effect_class(tool_name)
    if effect_class != "irreversible":
        return raw_executor(args, "") # Read-only tools pass through

    effective_thread = thread_id or self.current_thread_id or f"session-{args.get('user', self.default_thread_id)}"

    # 1. Compute Deterministic SHA-256 Idempotency Key
    idempotency_key = compute_idempotency_key(effective_thread, tool_name, args)

    # 2. Check Idempotency Cache (Suppress duplicate writes within thread)
    existing_ticket = self.lookup_committed_ticket(idempotency_key)
    if existing_ticket:
        logging.info(f"Idempotency Cache Hit: Key {idempotency_key[:10]}... Suppressing duplicate write.")
        return f"A refund ticket has been issued successfully with the amount of {int(existing_ticket['amount'])}"

    # 3. Policy Precondition Validation (Policy-as-Code)
    allowed, reason = policy_engine.validate_preconditions(tool_name, args)
    if not allowed:
        logging.warning(f"EffectGuard Gateway BLOCKED '{tool_name}': {reason}")
        return f"EffectGuard Gateway Error: {reason}"

    # 4. Write-Ahead Journaling: Status = PENDING
    journal_id = f"JRN-{uuid4().hex[:8].upper()}"
    self.record_journal_entry(journal_id, effective_thread, idempotency_key, tool_name, args, status="PENDING")

    # 5. Execution with Verify-Before-Retry
    try:
        result_str = raw_executor(args, idempotency_key)
        self.record_journal_entry(journal_id, effective_thread, idempotency_key, tool_name, args, status="COMMITTED", result_text=result_str)
        return result_str

    except TimeoutError as timeout_ex:
        logging.warning("Caught timeout during tool execution. Initiating Verify-Before-Retry...")
        committed_ticket = self.lookup_committed_ticket(idempotency_key)
        if committed_ticket:
            recovered_msg = f"A refund ticket has been issued successfully with the amount of {int(committed_ticket['amount'])}"
            self.record_journal_entry(journal_id, effective_thread, idempotency_key, tool_name, args, status="COMMITTED", result_text=recovered_msg)
            return recovered_msg

        # If not committed, retry once safely with identical key
        retry_result = raw_executor(args, idempotency_key)
        self.record_journal_entry(journal_id, effective_thread, idempotency_key, tool_name, args, status="COMMITTED", result_text=retry_result)
        return retry_result
\end{lstlisting}

\section{Policy Engine Implementation and Invariant Enforcement}
\label{sec:impl_policy}
The Policy Engine (\texttt{effectguard/policy\_engine.py}) parses \texttt{policy.yaml} and evaluates invariants deterministically in pure Python without requiring intermediate LLM calls.

The engine verifies the four primary preconditions:
\begin{enumerate}
    \item \textbf{User Existence:} Invokes \texttt{user\_exists(user)} against the SQLite \texttt{users} table.
    \item \textbf{Positive Amount:} Validates that $\text{float}(A) > 0$.
    \item \textbf{Safety Ceiling:} Enforces $A \le 500.0$.
    \item \textbf{Unrefunded Balance Verification:} Retrieves all purchases for user $u$ where $\texttt{status} = \texttt{'PURCHASED'}$, filters items strictly within the 30-day epoch window ($t_{\text{epoch}} - 30\text{ days} \le t_{\text{purchase}} \le t_{\text{epoch}}$), sums their prices, and guarantees that $A \le \text{MaxEligible}$. If $A$ exceeds this sum, execution is blocked.
\end{enumerate}

\section{Distributed Chaos Fault Injection Engine}
\label{sec:impl_chaos}
To simulate real-world distributed network failures, \texttt{agent/fault\_injector.py} intercepts tool calls using pre- and post-execution hooks. The chaos engine supports three fault injection modes:
\begin{itemize}
    \item \texttt{drop\_before\_exec}: Simulates an upstream gateway timeout (HTTP 504) or network disconnect before the tool begins execution.
    \item \texttt{timeout\_after\_commit}: Simulates the critical post-commit failure: the underlying database inserts the refund ticket and marks items as refunded, but the network connection abruptly drops before the confirmation packet reaches the agent runtime.
    \item \texttt{server\_error}: Simulates an unrecoverable internal server crash (HTTP 500).
\end{itemize}

Fault triggering is governed by a pseudo-random probability generator:
\begin{equation}
    \text{FaultTriggered} \iff \text{UniformRandom}(0.0, 1.0) < r_{\text{fault}}
\end{equation}
where $r_{\text{fault}} \in \{0.20, 0.35\}$.

\section{Claim-vs-Ledger Audit Engine}
\label{sec:impl_claim_auditor}
Implemented in \texttt{effectguard/claim\_auditor.py}, the Claim Auditor executes a two-stage extraction algorithm:
\begin{enumerate}
    \item \textbf{Structured Extraction:} Inspects the \texttt{claimed\_effects} array from the agent's Pydantic response. If populated, it extracts the claimed action and monetary sum directly.
    \item \textbf{Regex Text Fallback:} If \texttt{claimed\_effects} is empty or unpopulated by the model, the auditor executes calibrated regular expression extractors across the natural-language text to identify monetary claims and confirmation tokens (e.g., \texttt{r"refund\textbackslash s+(?:amount\textbackslash s+)?(?:of\textbackslash s+)?\textbackslash \$?(\textbackslash d+)"}).
\end{enumerate}

The extracted values ($\mathcal{C}_{\text{tickets}}, \mathcal{C}_{\text{amount}}$) are diffed against the physical database records ($\mathcal{L}_{\text{tickets}}, \mathcal{L}_{\text{amount}}$) retrieved via \texttt{SELECT * FROM refund\_tickets WHERE username = ?}. Discrepancies are flagged and exported into the telemetry stream.

\section{Scientific Matrix Runner and Ablation Suite}
\label{sec:impl_matrix}
Implemented in \texttt{effectguard/condition\_runner.py}, the Matrix Runner conducts comparative ablation experiments across five distinct architectural conditions:
\begin{itemize}
    \item \textbf{C0 (Naive Baseline):} Gateway disabled (\texttt{enabled = False}), blind retries upon receiving a timeout.
    \item \textbf{C1 (Exponential Backoff):} Gateway disabled, retries executed with an exponential backoff sleep delay ($t_{\text{sleep}} = 0.05 \times 2^i$ seconds).
    \item \textbf{C2 (Idempotency Keys Only):} Gateway enabled with SHA-256 keying, but Verify-Before-Retry disabled.
    \item \textbf{C3 (Verify-Before-Retry Only):} Database verification enabled on timeout, but idempotency keys replaced with random UUIDs.
    \item \textbf{C4 (Full EffectGuard):} All architectural safety mechanisms active (PAC + SHA-256 Keys + WAL Journal + Verify-Before-Retry).
\end{itemize}

\section{Relational Persistence and Concurrency Hardening}
\label{sec:impl_db}
In \texttt{agent/db.py}, the persistence layer is hardened for production concurrency:
\begin{lstlisting}[language=Python]
def get_db_connection() -> sqlite3.Connection:
    conn = sqlite3.connect(str(DB_PATH), timeout=20.0)
    conn.row_factory = sqlite3.Row
    conn.execute("PRAGMA journal_mode = WAL;")
    conn.execute("PRAGMA synchronous = NORMAL;")
    return conn
\end{lstlisting}
WAL mode ensures that read queries from the agent never block write transactions from the gateway, and write transactions never block concurrent evaluations.

\subsection{Username Normalization Algorithm}
To prevent capitalization and punctuation mismatches from bypassing policies, \texttt{normalize\_username(username)} implements a three-stage resolution:
\begin{enumerate}
    \item Direct case-insensitive lookup in \texttt{users} (\texttt{COLLATE NOCASE}).
    \item Underscore removal match (resolving \texttt{user\_123} $\leftrightarrow$ \texttt{user123}).
    \item Canonical fallback preserving valid alphanumeric handles (e.g., \texttt{sarah\_m}).
\end{enumerate}

\section{Hermetic 25-Scenario Benchmark Harness}
\label{sec:impl_benchmark_harness}
The benchmark harness (\texttt{evaluations/dataset.py} and \texttt{backend/eval\_service.py}) encapsulates 25 golden test scenarios categorized across five enterprise operational domains:
\begin{enumerate}
    \item \textbf{Valid Refunds (8 Scenarios):} Single-item and multi-item returns within the valid 30-day window (Cases 1--7, 9).
    \item \textbf{Policy Boundary Scenarios (4 Scenarios):} Edge-of-boundary testing including exactly 30 days old, 31 days old (expired), and invalid future purchase dates (Cases 8, 10, 11, 25).
    \item \textbf{Partial Eligibility Scenarios (5 Scenarios):} Multi-item returns containing mixed eligible, expired, and unpurchased items (Cases 12--16).
    \item \textbf{Informational Inquiries (4 Scenarios):} Inquiries regarding purchase history where no refund is requested and zero tickets must be issued (Cases 17--20).
    \item \textbf{Negative Guardrail Scenarios (4 Scenarios):} Unregistered users and unpurchased item return attempts (Cases 21--24).
\end{enumerate}

\subsection{Sequence Matching via Longest Common Subsequence (LCS)}
To verify that the agent invoked tools in the correct logical sequence (e.g., verifying user history \textit{before} calculating refunds, and calculating refunds \textit{before} issuing tickets), the evaluation harness implements an LCS-based sequence correctness metric:
\begin{equation}
    \text{Score}_{\text{seq}} = \frac{|\text{LCS}(\mathcal{T}_{\text{called}}, \mathcal{T}_{\text{expected}})|}{|\mathcal{T}_{\text{expected}}|}
\end{equation}

\section{Automated Testing Suite and Verification Matrix}
\label{sec:impl_testing}
To ensure rigorous software verification, the project maintains an automated test suite implemented in \texttt{pytest}. Table~\ref{tab:test_suite_matrix} summarizes the test modules, scope, and verification status across all 42 hermetic test cases.

\begin{table}[H]
\centering
\small
\renewcommand{\arraystretch}{1.2}
\begin{tabularx}{\textwidth}{p{4.5cm}cp{5.5cm}c}
\toprule
\textbf{Test Module Path} & \textbf{Tests} & \textbf{Verification Scope \& Assertions} & \textbf{Status} \\
\midrule
\texttt{tests/test\_frozen\_clock.py} & 3 & Verifies temporal epoch freezing, Day 29/30/31 boundary math. & \textbf{PASSED (3/3)} \\
\texttt{tests/test\_eval\_metrics.py} & 6 & Verifies LCS sequence matching, argument accuracy, financial metrics. & \textbf{PASSED (6/6)} \\
\texttt{tests/test\_phase1\_database...} & 6 & Verifies SQLite WAL writes, chaos injection, C0 duplicate baseline. & \textbf{PASSED (6/6)} \\
\texttt{tests/test\_phase2\_effectguard...} & 6 & Verifies Policy-as-Code preconditions, SHA-256 keys, Verify-Before-Retry. & \textbf{PASSED (6/6)} \\
\texttt{tests/test\_phase3\_adversarial...} & 7 & Verifies Claim Auditor diffs, adversarial prompt attacks, crash recovery. & \textbf{PASSED (7/7)} \\
\texttt{tests/test\_phase4\_matrix.py} & 2 & Verifies C0–C4 comparative matrix runner under chaos. & \textbf{PASSED (2/2)} \\
\texttt{tests/test\_persistence...} & 6 & Verifies item status updates (\texttt{REFUNDED}), username normalization. & \textbf{PASSED (6/6)} \\
\texttt{tests/test\_partial\_eligibility...}& 6 & Verifies multi-item cart staging and split refund calculations. & \textbf{PASSED (6/6)} \\
\midrule
\textbf{Total Automated Verification} & \textbf{42} & \textbf{Complete hermetic test coverage across all architectural components} & \textbf{100\% PASS} \\
\bottomrule
\end{tabularx}
\caption{Pytest test suite verification matrix across all 42 automated test cases.}
\label{tab:test_suite_matrix}
\end{table}
All 42 tests execute hermetically in approximately 2.0 seconds, confirming that all components operate in strict accordance with functional specifications.
